% ====================================================================================================
% LIMITATIONS + CONTRIBUTIONS  (unnumbered; not counted toward the 8 pages)
%   TODO(team): contribution statement and author names.
% ====================================================================================================

\section*{Limitations}
English only, four segmenters and UD annotation. The layout condition takes paragraph breaks from UD, where they are
sentence boundaries by construction, so results with paragraphs are optimistic (``structure given''). The held-out sets are
small (GENTLE 26 documents, LinES 9 and 18) and the new splits are new documents, not new domains. Baseline comparisons
and the sensitivity test are post hoc; only the final evaluator's scores on EWT dev and LinES train+dev are one-shot tests,
and its numbers on GENTLE, LinES test and EWT test are re-used tests. The training objective omits the document-end term
used in evaluation (harmless on GUM, but not unified). Training data overlap: SaT-sm and Stanza were trained on GUM (and
Stanza on EWT); our evidence and SaT-12l-sm are both supervised on GUM, and GENTLE follows GUM's conventions. The design
was adaptive (13 dated, pre-registered amendments); the document-end correction was found post hoc on a re-used test and
then confirmed on a new split. An independent code review found five issues that affected reported numbers; after the
fixes, the pre-registered selections and all one-shot hypotheses are unchanged (App.~\ref{app:protocol}). Selecting a
segmenter was not shown to beat always using SaT, and transfer to different annotation conventions (LinES) fails.

\section*{Contributions and Availability}
% TODO(team): confirm the team-size / contribution statement with the instructor.
Weifeng: method, experiments, analysis, \S3--5 and appendix. Yanhan: \S1--2, figures, examples, presentation. Both:
\S6--7. Code, configurations, Slurm scripts and all result files are in the project repository; every number in this
paper is produced by a script listed in App.~\ref{app:protocol}.
